\section{Data}\label{sec:data}

\subsection{Coseismic landslide inventories}
We selected nine earthquakes around the Pacific Rim that have openly available, event-wide landslide inventories (Figure~\ref{fig:events}; Table~\ref{tab:events}). The set spans subduction-interface, crustal-reverse and strike-slip ruptures, M\textsubscript{w} 6.1--7.9, and climates from sub-Arctic Alaska and Patagonian fjords to the humid tropics of New Guinea and Sulawesi. Eight inventories come from the USGS open repository of earthquake-triggered ground-failure inventories \citep{schmitt2017}: Northridge \citep{harp1996,data_harp}, Niigata-Chuetsu \citep{data_niigata}, Iwate-Miyagi \citep{yagi2009,data_iwate}, Ays\'en \citep{data_aisen}, Denali \citep{gorum2014denali,data_denali}, Cinchona \citep{data_cinchona}, Porgera \citep{tanyas2022png,data_porgera} and Palu \citep{zhao2021palu,data_palu}. For Kaik\={o}ura we used the centroids of the 29\,519 landslide source areas in version~2 of the GNS Science inventory \citep{massey2020update,data_kaikoura}, which were mapped from 0.3~m post-event aerial photographs and lidar. For Porgera, only landslides flagged as coseismic were kept; post-seismic landslides were excluded.

The inventories differ in format (polygons or points), in whether source and runout areas are separated, and in completeness \citep{tanyas2017}. To give all events a common sampling unit, we reduced each polygon to one interior representative point. Each event is therefore treated as a set of landslide locations at the resolution of the analysis grid (30~m), matching the point-based Kaik\={o}ura and Palu data. The study domain of each event is the published mapping extent where one exists (Niigata, Porgera and Palu). Otherwise it is the convex hull of the inventory buffered by 2~km, a conservative proxy for the area the mappers inspected. Sea and no-data cells were removed. A 30~m cell is labelled positive if it contains at least one landslide point. Landslide prevalence (the fraction of positive cells) ranges over two orders of magnitude between events (Table~\ref{tab:events}), partly because of physical differences and partly because of differences in mapping. This heterogeneity is part of what a transferred model must cope with, and we return to it in Section~\ref{sec:disc}.

\subsection{Remotely sensed terrain}
Terrain predictors were derived from the Copernicus GLO-30 digital surface model \citep{copdem}. GLO-30 is produced from TanDEM-X interferometric SAR acquisitions made between 2011 and 2015. An independent assessment found it the most accurate of the global 1-arcsecond elevation models \citep{guth2021}, and it is freely available within minutes for any point on Earth, which makes it suitable for rapid response. Tiles were read from the AWS open-data bucket and resampled bilinearly to a 30~m UTM grid for each event. For events before about 2011 the DEM post-dates the earthquake, so it may contain some post-failure topography; this affects at most the largest landslides. We computed eight terrain variables that do not depend on absolute elevation and therefore transfer across mountain ranges: slope; maximum slope in a $5\times5$ window; local relief and topographic position index in a 510~m window; relative slope position (0 = valley floor, 1 = ridge crest) in a 1~km window; profile and plan curvature from second-order finite differences of the elevation surface; and the standard deviation of slope in a $3\times3$ window as a roughness measure.

\subsection{Ground motion}
Peak ground acceleration (PGA) and peak ground velocity (PGV) were taken from the USGS ShakeMap grid of each event \citep{wald1999,worden2018} and interpolated bilinearly to the 30~m grid. Where both existed, the most recent contributed ShakeMap version was preferred over the ShakeMap Atlas version; this applies to Kaik\={o}ura, Porgera and Palu. ShakeMap quality varies strongly between events with station density. Sparse strong-motion networks leave the shaking field between stations uncertain, and adding GNSS-derived velocities can change the estimated PGV pattern enough to improve its agreement with landslide locations \citep{endo2026realtime}. For example, the Ays\'en and Cinchona ShakeMaps give maximum PGA of only 0.12~g and 0.27~g inside the landslide domains (Table~\ref{tab:events}), well below the shaking normally associated with dense landsliding \citep{keefer1984}. This variation is a realistic source of covariate shift \citep{shimodaira2000}, and we deliberately do not correct it, because an operational model would face the same uncertainty.

\subsection{Event-to-event differences in the landslide response}\label{sec:eventdiff}
Two exploratory views of the data motivate the design of the proposed algorithm. First, the predictor distributions differ strongly between events (Supplementary Figure~S1). Median PGV in the landslide domains ranges from 4~cm\,s$^{-1}$ (Ays\'en) and 10~cm\,s$^{-1}$ (Cinchona) to 55~cm\,s$^{-1}$ (Iwate). At Cinchona the PGV distributions of landslide and non-landslide cells overlap largely, so the ShakeMap carries little discriminating information inside that domain. A model transferred to these events must therefore extrapolate in its ground-motion inputs, which is a direct form of covariate shift \citep{shimodaira2000}.

Second, the empirical frequency of landsliding, binned by Newmark displacement or by PGA, follows broadly similar increasing curves in all events, but the curves are displaced vertically by up to two orders of magnitude (Figure~\ref{fig:response}). For cells with a computed Newmark displacement of about 1~cm, the landslide frequency is 0.04\% at Denali, 0.44\% at Kaik\={o}ura and 2.3--2.6\% at Niigata-Chuetsu and Cinchona. Within an event, the variation of landslide frequency with displacement is about as large as this between-event spread. Much of the difficulty of transfer is therefore a difference in \emph{level} rather than in \emph{shape}. A model that shares the shape across events and lets the level vary by event, which is what HiPhyBoost does through its offsets $b_e$, fits this structure; a single pooled response curve cannot.

\begin{figure*}[t]\centering
\includegraphics[width=\textwidth]{fig_response_curves.pdf}
\caption{Empirical landslide frequency (weighted fraction of 30~m cells containing a landslide) as a function of (a) Newmark displacement computed with uniform strength parameters and (b) ShakeMap PGA, for the nine earthquakes. Bins with fewer than five landslides are omitted. The curves have similar increasing shapes but are offset vertically by up to two orders of magnitude, which motivates the latent event offsets of HiPhyBoost.}\label{fig:response}
\end{figure*}
